\subsection{Los dos funcionales del perfil}\label{sec:profile-data}
Escribamos \(z=(r,q)\in[0,1]^2\). La densidad límite \(\nu(z)\) tiene marginales uniformes en fila y
columna. Las coordenadas diagonales son \(d=r-q\) y \(a=r+q\); en sus fibras se integra en \(dr\).
Si \(m_f(y)\) es la integral de \(\nu\) sobre la línea de coordenada \(y\) en la familia \(f=D,A\),
su vacancia final es \(u_f(y)=1-m_f(y)\). Para la fracción residual \(x\) ponemos
\[
 U_f(y;x)=u_f(y)+x(1-u_f(y)),\qquad U_R=U_C=x,\qquad
 W(z)=\frac{\nu(z)}{U_D(d;x)U_A(a;x)} .
\]
\paragraph{Corrección de torres.}
Escalamos \(v_x(z)=U_D(d;x)U_A(a;x)\) a marginales de masa uno,
\(p_x(r,q)=s_x(r)v_x(r,q)s_x(q)\); la simetría da el mismo factor a ambos lados.
Multiplicar \(v_x\) por el factor separable \(\nu/(u_Du_A)\) da el mismo núcleo escalado.
En \(x=1\) el núcleo escalado vale uno. Sea \(B(r,q)=p_x(r,q)-1\) como núcleo en \(L^2[0,1]\).
Los productos de operadores se definen por integración. Ponemos
\[
 T=(I-B^2)^{-1}-I,\qquad H=TB,\qquad
 j(r)=(T-B^2)(r,r),\qquad j_0=\int_0^1j(r)\,dr .
\]
Aquí \(\langle f\rangle=\int_{[0,1]^2} f(r,q)\,dr\,dq\).
Las potencias dentro de estos promedios son potencias puntuales de núcleos.
\begin{align*}
 F(x)={}&\tfrac12 j_0^2+\operatorname{tr}(T^2)-\tfrac32j_0-\tfrac32\int_0^1 j(r)^2\,dr\\
 &+\langle j(r)(B+H-T)(r,q)j(q)\rangle\\
 &-\tfrac23\langle T^3\rangle+2\langle B H^2\rangle+\tfrac23\langle H^3\rangle .
\end{align*}
En particular \(B=0\) para \(x=1\), lo que da \(F(1)=0\).
\paragraph{Prefactor de líneas.}
Para \(x>0\) usemos una malla de puntos medios de lado par \(m\); sean \(J(z)\) las cuatro líneas
que pasan por la casilla \(z\). Hay \(m,m,2m-1,2m-1\) líneas de tipos \(R,C,D,A\). Ponemos
\(\rho_l=1/m\) en filas y columnas, y \(\rho_l=U_l/(xm)\) en diagonales de ambas familias.
En el cruce \(z\) de dos líneas distintas ponemos
\[
 K_{lh}=W(z)\prod_{a\in J(z)\setminus\{l,h\}}U_a,\qquad
 v_R=1,\quad v_C=-1,\quad v_D=v_A=0,\qquad \Pi=vv^{\mathsf T}/2.
\]
Las demás entradas de \(K\) valen cero. Una diagonal y una antidiagonal se cruzan en la malla
sólo cuando sus paridades coinciden; así surgen las dos hojas de cada familia diagonal en el continuo.
Con \(\rho\) también para la matriz diagonal, definimos
\begin{align*}
 L_0&=K(I+\rho K+\rho\Pi)^{-1}\rho K-3\Pi,\qquad
 \eta_l=\begin{cases}1,&l\in R,C,\\ x/U_l,&l\in D,A,\end{cases}\\
 s_l&=m^{-1}\sum_{z\in l}\nu(z)W(z),\qquad
 \theta_l=3\eta_l s_l-\sum_h K_{lh}^2\rho_h .
\end{align*}
Sea \([f]_m=m^{-2}\sum_z\nu(z)f(z)\). El prefactor de líneas del continuo es el límite de
\begin{align*}
 \mathcal V_m(x)={}&4[W^2]_m
 -\tfrac32\Big[W\sum_{l,h\in J(z)}(L_0)_{lh}\Big]_m\\
 &-\tfrac32\Big[W\sum_{l\in J(z)}\theta_l\Big]_m
 +\frac{9}{4m}\sum_l\eta_l s_l^2 .
\end{align*}
La suma sobre \(l,h\) es ordenada e incluye su diagonal. El valor en cero se obtiene por
continuación desde \(x>0\). El Apéndice~\ref{app:interior} identifica los términos de
inserción de \(d(x)\) con este mismo perfil y la misma convención de paridad.
